%!TEX root = iopjournal.tex
%
% Related Work
%
% Outline (four parts, settled 04.09.2026). The TEP comes last so that the
% funnel narrows monotonically -- model class, application studies, requirement
% on the data basis, the one candidate that meets it, what it still lacks --
% and so that sections 2 and 3 run in the same order, concept drift first and
% the TEP last, with section 3 handing the plant straight to section 4.
%
% Each part closes with the residue it leaves open. The sum of the residues is
% the gap that IDV(29) fills.
%
%   2.1 Adaptive soft sensors in the process industry
%       - the model class the cost surrogate belongs to; the established
%         adaptation mechanisms (moving window, recursive schemes, just-in-time
%         learning, ensembles) and their precondition, a label available online
%         [kadlec.2009, kadlec.2011]
%       - TODO a recent survey on adaptive soft sensors (2020 onwards); in
%         neither references.bib nor diss.bib
%       - Residue: the causes listed there are all on the process side, the
%         measurement chain does not appear among them
%
%   2.2 Concept drift detection and adaptation in industrial applications
%       - studies, not methods: electroslag remelting [bayram.2022b] and the own
%         grinding plant [marijan.2025]; the application side is predominantly
%         classification [lu.2018, bayram.2022 as the frame]
%       - TODO two or three further industrial case studies
%       - Residue: onset and extent of the drift are unknown, so detectors can
%         only be compared with one another rather than measured
%
%   2.3 Benchmark streams with a known drift ground truth
%       - synthetic generators: exact drift onset, no plant physics
%       - real-world streams: physics, but no ground truth [souza.2020]
%       - regression streams are the exception in both camps [webb.2016]
%       - TODO generator and framework literature (MOA, river, SEA / rotating
%         hyperplane / Agrawal); in neither bib file
%       - Residue, phrased as a demand rather than a conclusion, so that 2.4
%         reads as the answer to it: what is required is a simulated plant that
%         carries plant physics and a defined drift onset at the same time
%
%   2.4 The Tennessee Eastman process as a benchmark
%       - history, model revisions and control concepts [downs.1993, ricker.1995,
%         ricker.1996, ricker.2005, bathelt.2015, reinartz.2022]
%       - what the TEP is actually used for: fault detection and diagnosis; the
%         standardised public reference data is cut for classification
%         [reinartz.2021]
%       - TODO three to five TEP fault-detection and -diagnosis works (PCA/DPCA,
%         deep learning). This is the largest of the four gaps, diss.bib holds
%         none of them
%       - the disturbance catalogue is characterised as a FAULT catalogue only;
%         IDV(13) at most in a subordinate clause as the single candidate
%         labelled "slow drift". The analysis stays in sec_drift_injection --
%         do not pre-empt the own finding here
%       - closing paragraph = the gap: the TEP meets every requirement derived
%         in the introduction except the drift itself
%
% Boundary to the background: related work reports STUDIES and DATA HOLDINGS,
% the background DEFINITIONS and MECHANICS. [gama.2004, baena.2006, bifet.2007,
% raab.2020, widmer.1996, moreno-torres.2012] belong to section 3, and so does
% the TEP as a plant -- components, control concept, cost function.
%
% Cf. dissertation: chapter 2 and sec:TEP
%

\section{Related Work} \label{sec:related_work}

\subsection{Adaptive soft sensors in the process industry} \label{subsec:soft_sensors}

\subsection{Concept drift detection and adaptation in industrial applications} \label{subsec:cd_industrial}

\subsection{Benchmark streams with a known drift ground truth} \label{subsec:benchmarks}

\subsection{The Tennessee Eastman process as a benchmark} \label{subsec:tep_benchmark}
